\section{Spatial discretization and physical time}
\label{sec:grid}

\subsection{Cubic lattice and voxelized sensor surface}

The explicit domain has nominal dimensions $L_x\times L_y\times H$ and is
partitioned into a simple cubic lattice with spacing $a$.  The implemented
numbers of lattice intervals are
\begin{equation}
 N_x=\operatorname{round}\!\left(\frac{L_x}{a}\right),\qquad
 N_y=\operatorname{round}\!\left(\frac{L_y}{a}\right),\qquad
 N_z=\operatorname{round}\!\left(\frac{H}{a}\right).
 \label{eq:grid-counts}
\end{equation}
Fluid and solid sites are stored on the resulting Cartesian grid.  An
analytical surface---planar, cylindrically curved, or sinusoidal---is converted
to a solid mask.  Every face separating an accessible fluid voxel from a
solid voxel is an exposed surface face with area $a^2$.  Consequently, the
voxelized reactive area is
\begin{equation}
 A_s=\sum_{f\in\mathcal F_{\mathrm{reactive}}}a^2,
 \label{eq:voxel-area}
\end{equation}
where $\mathcal F_{\mathrm{reactive}}$ contains the bulk-accessible exposed
faces.  At prescribed receptor density $\rho_R$, the receptor count is
\begin{equation}
 N_R=\operatorname{round}(\rho_R A_s),
 \label{eq:receptor-count}
\end{equation}
unless a fixed receptor count is explicitly imposed.  Receptor-bearing faces
are sampled uniformly without replacement for the manuscript simulations;
there is therefore at most one receptor per surface face.  Fixed-$N_R$
controls separate geometric confinement from changes in total receptor
number.

The staircase representation in Eq.~\ref{eq:voxel-area} is part of the
numerical model.  Refining $a$ improves geometric resolution but also changes
the local reaction volume and, in a local lattice reaction model, can change
the relationship between microscopic and mesoscopic reaction rates.  Spatial
resolution must therefore be assessed by convergence studies rather than by
assuming that arbitrarily small voxels are automatically more physical
\citep{isaacson2008,hellander2012,hellander2015}.

\subsection{Diffusive update interval}

The coefficient $\lambda_D$ is not an additional physical parameter; it is
the dimensionless group produced by discretizing Fick's second law.  Let
$c_{i,j,k}^{n}$ denote the concentration at lattice site $(i,j,k)$ and time
$t_n=n\dt$.  A forward difference in time and centered second differences in
space give
\begin{align}
 \frac{c_{i,j,k}^{n+1}-c_{i,j,k}^{n}}{\dt}
 =\frac{D}{a^2}\big(&c_{i+1,j,k}^{n}+c_{i-1,j,k}^{n}
 +c_{i,j+1,k}^{n}+c_{i,j-1,k}^{n}\nonumber\\
 &+c_{i,j,k+1}^{n}+c_{i,j,k-1}^{n}
 -6c_{i,j,k}^{n}\big).
 \label{eq:ftcs-grid-derivation}
\end{align}
Multiplying Eq.~\ref{eq:ftcs-grid-derivation} by $\dt$ shows that the fraction
transferred from a site to \emph{each} nearest neighbor during one update is
$D\dt/a^2$.  This is the standard forward-time, centered-space (FTCS)
discretization of the diffusion equation \citep{crank1975}.  We therefore
define
\begin{equation}
 \lambda_D=\frac{D\dt}{a^2}.
 \label{eq:lambda-D}
\end{equation}
It can also be read as the ratio between the update duration $\dt$ and the
diffusive timescale $a^2/D$ associated with one lattice spacing.  In the
single-particle representation, the ligand moves to any of its six nearest
neighbors with probability $\lambda_D$ per direction and remains at its
present site with probability
\begin{equation}
 p_0=1-6\lambda_D.
 \label{eq:stay-probability}
\end{equation}
Nonnegative probabilities require the standard three-dimensional explicit
diffusion stability condition
\begin{equation}
 \dt\leq \frac{a^2}{6D}.
 \label{eq:stability}
\end{equation}
When the user does not prescribe a timestep, the implementation uses
\begin{equation}
 \dt=0.95\frac{a^2}{6D},
 \label{eq:default-dt}
\end{equation}
which gives $p_0=0.05$ and $p_{\pm x}=p_{\pm y}=p_{\pm z}=0.95/6$.  For the
manuscript baseline $a=\SI{1}{\nano\meter}$ and
$D=\SI{1e-11}{\meter\squared\per\second}$,
Eq.~\ref{eq:default-dt} gives $\dt\approx\SI{1.58e-8}{\second}$.  Physical
observation times are obtained as integer numbers of these updates.

\subsection{Numerical scales that must remain distinct}

Three lengths have different meanings and should not be conflated:
\begin{enumerate}[leftmargin=*]
 \item $a$ controls the spatial resolution, surface voxelization, and local
 reaction volume;
 \item the geometric length scales (for example, a cylindrical radius or
 sinusoidal wavelength) define confinement and curvature; and
 \item the escape height defines the boundary of the explicitly simulated
 near-surface region and hence the first-passage endpoint.
\end{enumerate}
The supplementary numerical controls vary $a$, lateral domain size, and
escape height because each can otherwise introduce an arbitrary computational
scale into a rebinding calculation.
